\documentclass{article}
\usepackage{mathnotes}

\title{Euler's Method}
\description{Euler's numerical methods for approximating solutions to initial value problems for ordinary differential equations.}

\begin{document}

\section{Euler's Method}

We'll consider a first order IVP of the form

\[ \frac{du}{dt} f(t, u), \quad u(0) = c, \]

with the goal of numerically estimating the solution on a finite time interval $[0, T]$ for some $T > 0.$

First we'll establish this is a reasonable thing to do.

\begin{theorem}[Well-posedness of IVPs]
Consider the IVP

	\[ \frac{d\vec{u}}{dt} = \vec{f}(t, \vec{u}), \quad \vec{u}(0) = \vec{c}. \]

	Suppose that $\vec{f}: [0, T] \times \reals^d \to \reals^d$ is \@{continuous} in its \@{domain} and \@{Lipschitz continuous} in $\vec{u}$, i.e., there exists a constant $L > 0$ such that

	\[ |\vec{f}(t, \vec{u}) - \vec{f}(t, \vec{v})| \leq L\|\vec{u} - \vec{v}\|, \forall(t, \vec{u}), (t, \vec{v}) \in [0, T] \times \reals^d. \]

	Then it holds that:

	\begin{enumerate}
		\item \textbf{Existence and Uniqueness:} There exists a unique solution $\vec{u}(t)$ for any $\vec{u}_0 \in \reals^d.$
		\item \textbf{Continuous Dependence on initial data and right hand side:} Suppose $\vec{r} : [0,T] \times \reals^d \to \reals^d$ is a bounded \@{perturbation} of $\vec{f},$ i.e. there exists a constant $M > 0$ such that
			\[ \| \vec{r}(t, \vec{u}) \| \leq M, \quad \forall(t, \vec{u}) \in [0, T] \times \reals^d. \]
			Further let $\vec{\hat{u}}(t)$ be the solution of the perturbed IVP

			\[ \frac{d \vec{\hat{u}}}{dt} = \vec{f}(t, \vec{\hat{u}}) + \vec{r}(t, \vec{\hat{u}}), \quad \vec{\hat{u}}(0) = \vec{\hat{c}},  \]
			which we assume exists and is unique. Then we have

			\[ \|\vec{u}(t) - \vec{\hat{u}}(t) \| \leq e^{Lt} \| \vec{c} - \vec{\hat{c}} \| + \frac{M}{L}(e^{Lt} - 1), \quad \forall t \in [0, T]. \]
	\end{enumerate}

\end{theorem}

Now that we've established that IVPs are well posed, we can look at how to approximate our unknown function $u(t)$ at a discrete set of points in time. 

\begin{definition}[discrete time points]\synonyms{"time mesh points", "time grid points"}
	Consider a \@{partition} $0 = t_0 < t_1 < \cdots < t_N = T$ of the time interval $[0, T].$ We can call the points $t_n$ the \textbf{discrete time points}. Then our goal is to approximate the values $u(t_n)$ with a set of quantities $\{u_n\}_{n=0}^N,$ that is, $u_n \approx u(t_n)$ for $n = 0, 1, \dots, N.$ We can also write $\delta t_n = t_{n+1} - t_n$ to denote the \term[time step]{step size} between between consecutive points. For simplicity we often assume that $\delta t_n = \delta t$ for all $n = 0, 1, \dots, N-1,$ i.e. the time steps are constant. Such a mesh is called a \@{uniform mesh}.
\end{definition}

\end{document}
